\subsection{Realization over prime fields}

The rank-two bound is sharp for prime-field data even when the
auxiliary torsion level is fixed arbitrarily high.

\begin{proposition}[Finite-field sharpness]\label{prop:sharpness}
Fix a totally real field \(F\) of degree \(e\), put \(E=F(i)\),
and fix an odd rational prime \(\ell\) and \(N\ge1\). For each
\(0\le r<e/2\), and for every sufficiently large rational prime
\(p\) split completely in \(E\) and congruent to \(1\) modulo
\(\ell^{2N}\), there is an abelian variety over \(\mathbb F_p\)
of dimension \(2e\) with the following properties. It admits a
Rosati-compatible rank-two action by \(E\), all \(e\) split pairs
have dimensions \((1,1)\), exactly \(r\) pairs are nonordinary,
and its \(\ell^N\)-torsion is rational.
\end{proposition}

\begin{proof}
Let \(M=\ell^{2N}\) and let \(I\subset\mathcal O_F\) be the product
of any \(r\) distinct primes over \(p\). We first choose
\(\tau\in I\) with \(\tau\equiv 2\pmod{M\mathcal O_F}\) and
\(\Q(\tau)=F\), such that \(|\sigma(\tau)|\ll_{F,M}p^{r/e}\)
for each real embedding \(\sigma\).

Here is the lattice argument. Write \(I=\beta J\), where \(J\) comes
from a fixed finite list of fractional ideal class representatives.
The logarithmic unit lattice of \(F\) has a bounded fundamental
domain in the hyperplane of coordinate sum zero. Multiplying
\(\beta\) by a unit lets us bound each \(|\sigma(\beta)|\) by
a constant times \(|\Nm_{F/\Q}(\beta)|^{1/e}\), hence by a constant
times \(p^{r/e}\). Transport a fixed \(\Z\)-basis of the lattice
\(J\) to a basis of \(I\). All its coordinates in the real embeddings
are bounded by \(O_F(p^{r/e})\). Since \(I\) is coprime to \(M\), the
congruence is a nonempty coset modulo \(MI\) in \(I\). Reducing a
representative in a fundamental parallelepiped gives the required
size bound. If \(F\) has \(s\) proper subfields, vary this
representative by \(MI\)-basis combinations with coefficients in
\(\{0,\ldots,s\}\). In grid coordinates each proper subfield
meets at most \((s+1)^{e-1}\) points, since it lies in a
proper affine hyperplane. Some point avoids all the subfields.
This preserves the bound with a larger fixed constant.

For sufficiently large \(p\), the strict inequality
\(|\sigma(\tau)|<2\sqrt p\) follows from \(r/e<1/2\).
The norm is nonzero and bounded in absolute value by
\(O_{F,M}(p^r)\), so \(p^{r+1}\) cannot divide it. Thus
\(v_{\mathfrak q}(\tau)=1\) at the selected primes, and
\(\tau\) is a unit at the other primes of \(F\) over \(p\).
Let \(\pi\) be a root of \(X^2-\tau X+p\). This is a Weil
\(p\)-integer, since the quadratics in all real embeddings
have nonreal conjugate roots of modulus \(\sqrt p\).
The field \(C=\Q(\pi)=F(\pi)\) is CM of degree \(2e\),
using \(\tau=\pi+p/\pi\). At the selected primes the
quadratic is Eisenstein with slopes \((1/2,1/2)\).
At every other prime over \(p\) it splits by Hensel's lemma
with slopes \((0,1)\).

By Honda--Tate \cite[Theorem~1]{tate}, \(\pi\) is the Frobenius
of a simple abelian variety \(B/\mathbb F_p\) up to isogeny.
The only potentially fractional local endomorphism invariant
over \(p\) has the form \(2\cdot(1/2)=1\) modulo \(1\).
All invariants vanish, so the rational endomorphism algebra over
\(\mathbb F_p\) is \(C\), and \(B\) has dimension \(e\).
By Waterhouse \cite[Theorem~6.1]{waterhouse}, choose \(B\) with
endomorphism ring \(\mathcal O_C\).
The number \(\eta=(\pi-1)/\ell^N\) belongs to that ring, since
it satisfies the monic equation
\[
 \eta^2+\frac{2-\tau}{\ell^N}\eta
       +\frac{1-\tau+p}{\ell^{2N}}=0
\]
over \(\mathcal O_F\). Therefore Frobenius acts trivially on
\(B[\ell^N]\).

On \(A=B^2\), let \(F\) act diagonally and let \(i\) act by
the matrix with rows \((0,-1)\) and \((1,0)\).
On \(B\), Rosati sends \(\pi\) to \(p/\pi\), so it fixes
\(F=\Q(\pi+p/\pi)\). The product Rosati involution fixes the
diagonal action and sends the rotation matrix to its negative.
Over \(\Q_p\), the height-two \(F\)-components of \(B\)
follow from the relative quadratic Frobenius polynomial.
The action of \(i\) splits at both places above each prime of \(F\).
Each of the two eigenspaces inherits the height-two component
of \(B\) at that prime, with slope sum one. Thus exactly
the selected \(r\) pairs are nonordinary, and all \(\ell^N\)-torsion
is rational as claimed.
\end{proof}

The primes in the proposition exist by Chebotarev. Taking
\(r=\lfloor(e-1)/2\rfloor\) realizes the bound for every all-balanced
signature over prime fields. The examples vary with \(p\).
A conclusion about the density of the residual pattern for a fixed
characteristic-zero variety requires global input beyond these
split Frobenius and torsion data.
